% ===========================================================================
%  Capítulo 2 — Análisis del dominio de planificación tabular
% ===========================================================================
\chapter{An\'alisis del dominio de planificaci\'on tabular}
\label{cap:dominio}

Para dise\~nar un lenguaje que capture los patrones recurrentes de la
planificaci\'on tabular, es necesario estudiar primero los casos concretos
que el lenguaje debe resolver.  En este cap\'{\i}tulo se presentan tres
casos de uso reales de la Facultad de Matem\'atica y Computaci\'on de la
Universidad de La Habana, se identifican los elementos comunes entre ellos
y se derivan los requisitos del lenguaje.

% ---------------------------------------------------------------------------
\section{Casos de uso}
\label{sec:casos}
% ---------------------------------------------------------------------------

% ............................................................................
\subsection{Horario de defensas de tesis}
\label{ssec:caso-defensas}
% ............................................................................

El proceso de planificaci\'on de defensas de tesis de diploma asigna a cada
estudiante un conjunto de cinco profesores que forman el tribunal ---tutor,
oponente, presidente, vocal y secretario---, una fecha y hora de exposici\'on,
y un local donde esta se celebra.  El coordinador de defensas debe detectar
conflictos: un mismo profesor asignado a dos defensas con el mismo slot
horario (mismo d\'{\i}a y hora), lo que hace imposible su presencia en
ambas.

La hoja de c\'alculo que da soporte a este proceso contiene dos tablas:

\begin{description}
  \item[Tabla de defensas.] Una fila por defensa, con columnas para el
    estudiante y los cinco roles del tribunal, m\'as las columnas de d\'{\i}a,
    hora y local.  Esta tabla no necesita columnas calculadas: toda la
    informaci\'on se introduce manualmente.

  \item[Tabla de profesores.] Una fila por profesor, con su nombre, grado
    acad\'emico y el n\'umero de defensas en que participa.  Este \'ultimo dato
    se calcula autom\'aticamente contando cu\'antas veces aparece el nombre del
    profesor en cualquiera de los cinco roles de la tabla de defensas.
    Adem\'as, las filas de los profesores con conflicto de horario se resaltan
    en rojo.
\end{description}

Los datos del caso de validaci\'on incluyen tres defensas y seis profesores.
Dos de las defensas comparten el slot \texttt{Lunes~10:00} con integrantes
de tribunal en com\'un, lo que genera un conflicto que la hoja debe detectar
y resaltar:

\begin{table}[h!]
\centering
\caption{Datos de las defensas del caso de validaci\'on.}
\label{tab:datos-defensas}
\renewcommand{\arraystretch}{1.2}
\begin{tabular}{llllllll}
\hline
\textbf{Estudiante} & \textbf{Tutor} & \textbf{Opon.}
  & \textbf{Pres.} & \textbf{Vocal} & \textbf{Secr.}
  & \textbf{D\'{\i}a} & \textbf{Hora} \\
\hline
Juan D\'{\i}az  & Piad   & Garc\'{\i}a & L\'opez  & Torres & Ruiz   & Lunes & 10:00 \\
Mar\'{\i}a Vega & Garc\'{\i}a & L\'opez  & Piad   & Soto   & Torres & Lunes & 14:00 \\
Luis Mora   & L\'opez  & Torres & Garc\'{\i}a & Ruiz   & Soto   & Lunes & 10:00 \\
\hline
\end{tabular}
\end{table}

Las defensas de Juan D\'{\i}az y Luis Mora comparten el slot
\texttt{Lunes~10:00}.  El profesor L\'opez participa como oponente en la
primera y como tutor en la segunda; el profesor Torres como vocal en la
primera y como oponente en la segunda.  Ambos tienen conflicto y deben
aparecer resaltados en rojo en la tabla de profesores.

Este caso es el que se usa como ejemplo motivador a lo largo del cap\'{\i}tulo~\ref{cap:dsl}.

% ............................................................................
\subsection{Horario docente de la facultad}
\label{ssec:caso-facultad}
% ............................................................................

La asignaci\'on de docencia en la Facultad de Matem\'atica y Computaci\'on
implica distribuir los turnos de clase de cada semana entre los grupos de las
tres carreras (Ciencia de la Computaci\'on, Matem\'atica y Ciencia de Datos),
en cuatro a\~nos cada una.  Cada grupo tiene una hoja propia en el workbook
con su horario semanal.

La hoja de c\'alculo de planificaci\'on contiene, para cada grupo, tres tablas
dispuestas en la misma fila de la hoja:

\begin{description}
  \item[Tabla de turnos (\texttt{turno-table}).] Seis columnas (turno,
    lunes, martes, mi\'ercoles, jueves y viernes), con dos filas Excel por
    turno l\'ogico (\dsl{cell-height~=~2}): la primera subfila almacena el
    nombre de la asignatura y la segunda el aula asignada.

  \item[Tabla de estad\'{\i}sticas (\texttt{stats-table}).] Una fila por
    asignatura, con su abreviatura, nombre completo, frecuencia semanal
    requerida, asignaciones ya realizadas (calculadas contando las apariciones
    de la abreviatura en la tabla de turnos) y horas que faltan por asignar
    (calculadas como la diferencia entre frecuencia y asignadas).

  \item[Tabla de aulas (\texttt{aulas-table}).] Lista de aulas disponibles
    para el grupo.
\end{description}

Adem\'as de las hojas de grupo, existe una hoja de ocupaci\'on de aulas que
agrega, para cada turno de cada d\'{\i}a, los grupos que ocupan cada aula.
Esta tabla usa referencias cruzadas din\'amicas (\dsl{collect-over}) que
iteran sobre todas las hojas de grupo y extraen el aula asignada en cada
turno.  Diecisiete grupos generan diecisiete hojas y una hoja de aulas con
referencias a las diecisiete.

Este caso pone a prueba dos caracter\'{\i}sticas del DSL que el caso de
defensas no ejercita: las tablas con celdas de altura compuesta
(\dsl{cell-height} mayor que 1) y las referencias cruzadas entre hojas del
mismo workbook mediante \dsl{collect-over}.

% ............................................................................
\subsection{Horario de programaci\'on televisiva}
\label{ssec:caso-tv}
% ............................................................................

El tercer caso de uso es el horario de programaci\'on semanal de un canal de
televisi\'on.  Cada programa tiene una hora de inicio, una duraci\'on en
minutos, un nombre, un tipo y un p\'ublico objetivo.  La hoja de c\'alculo
muestra una p\'agina por d\'{\i}a, con columnas para cada uno de esos campos
m\'as una columna de hora de fin, calculada autom\'aticamente a partir de la
hora de inicio y la duraci\'on.

La f\'ormula de la hora de fin involucra operaciones de tiempo:
\[
\texttt{fin} = \texttt{TEXT}(\texttt{TIMEVALUE}(\texttt{inicio}) + \texttt{duraci\'on}/1440, \texttt{"hh:mm"})
\]
donde $1440$ es el n\'umero de minutos en un d\'{\i}a.  Adem\'as, cada fila
se colorea seg\'un el tipo de programa (informativo, musical, infantil, etc.)
mediante un color distinto en la celda.

Este caso valida que el DSL puede expresar operaciones aritm\'eticas sobre
valores de tiempo (\dsl{time-add}) y f\'ormulas de celda simples sin
cuantificadores existenciales.

% ---------------------------------------------------------------------------
\section{Elementos comunes}
\label{sec:elementos-comunes}
% ---------------------------------------------------------------------------

Del an\'alisis de los tres casos se identifican cuatro patrones que se
repiten en todos o en la mayor\'{\i}a de los documentos de planificaci\'on
tabular:

% ............................................................................
\subsection{Tablas relacionadas y referencias cruzadas}
\label{ssec:elem-tablas}
% ............................................................................

En todos los casos, el documento contiene m\'as de una tabla cuyos datos est\'an
relacionados.  La tabla de profesores refencia la tabla de defensas; la tabla
de estad\'{\i}sticas referencia la tabla de turnos; la hoja de aulas referencia
las hojas de cada grupo.

Estas referencias son de dos tipos:
\begin{itemize}
  \item \textbf{Intrarregi\'on:} ambas tablas est\'an en la misma hoja y
    se accede a ellas mediante rangos absolutos de columnas (p.~ej.,
    \texttt{\$B\$4:\$F\$6} para los cinco roles de la tabla de defensas).
  \item \textbf{Interhoja:} la tabla referenciada est\'a en otra hoja del
    mismo workbook (p.~ej., la hoja de aulas referencia las hojas de grupos).
\end{itemize}

El lenguaje debe resolver estas referencias sin que el programador escriba
letras de columna ni nombres de hoja: los identifica por los s\'{\i}mbolos
DSL de las columnas y por el identificador de la tabla.

% ............................................................................
\subsection{Columnas calculadas}
\label{ssec:elem-computed}
% ............................................................................

En todos los casos, al menos una columna de alguna tabla muestra un valor
calculado a partir de datos de otra tabla.  Los tipos de c\'alculo m\'as
frecuentes son:

\begin{description}
  \item[\texttt{COUNTIF}.] Cuenta cu\'antas celdas de un rango cumplen un
    criterio.  Ejemplo: n\'umero de veces que aparece el nombre de un
    profesor en los cinco roles de la tabla de defensas.

  \item[\texttt{COUNTA}.] Cuenta celdas no vac\'{\i}as en un rango.  Ejemplo:
    n\'umero de defensas registradas.

  \item[\texttt{SUM}.] Suma un rango.  Ejemplo: total de defensas acumuladas
    por todos los profesores.

  \item[Aritmética de tiempo.] Suma una duraci\'on en minutos a una hora de
    inicio.  Ejemplo: hora de fin de un programa televisivo.

  \item[\texttt{collect-over}.] Agrega resultados de la misma posici\'on en
    m\'ultiples hojas.  Ejemplo: grupos que ocupan el Aula~1 en el turno~3
    del lunes.
\end{description}

% ............................................................................
\subsection{Visualizaci\'on condicional}
\label{ssec:elem-condicional}
% ............................................................................

En los casos de defensas y de docencia, la condici\'on de coloreado depende
de datos de otra tabla.  El color no es est\'atico: si el usuario edita los
datos directamente en Excel, el coloreado debe recalcularse autom\'aticamente.

El mecanismo adecuado es, por tanto, una \texttt{FormulaRule} de formato
condicional incrustada en el archivo \xlsx{}, no un coloreado est\'atico
aplicado en tiempo de generaci\'on.  La condici\'on de la regla es una
f\'ormula SUMPRODUCT que captura la sem\'antica del cuantificador existencial
del DSL.

La forma en que el DSL expresa esta condici\'on ---mediante \dsl{exists}---
y la forma en que el backend la implementa ---mediante SUMPRODUCT--- se
detallan en los cap\'{\i}tulos~\ref{cap:dsl} y~\ref{cap:generacion},
respectivamente.

% ............................................................................
\subsection{Posicionamiento relativo}
\label{ssec:elem-posicion}
% ............................................................................

Todos los casos requieren expresiones fuera de las tablas: totales, etiquetas,
celdas de resumen.  La posici\'on de estas expresiones depende de cu\'antas
filas de datos tiene la tabla, que var\'{\i}a en cada ejecuci\'on.

La soluci\'on adoptada en el DSL es el sistema de anclas y navegaci\'on
(\dsl{nav}, \dsl{ultima-fila}, \dsl{primera-fila}): se declara que una
expresi\'on va ``una fila por debajo de la \'ultima fila de datos de la
columna \texttt{estudiante} de la tabla de defensas'', y el compilador
resuelve esa declaraci\'on a la coordenada Excel concreta en tiempo de
generaci\'on.

% ---------------------------------------------------------------------------
\section{Requisitos del lenguaje}
\label{sec:requisitos}
% ---------------------------------------------------------------------------

Del an\'alisis anterior se derivan los requisitos que el DSL debe satisfacer:

\begin{enumerate}
  \item \textbf{Declaraci\'on de tablas.} El lenguaje debe permitir definir
    tablas con nombre, con una lista de columnas identificadas por s\'{\i}mbolos,
    y con soporte para celdas de altura y anchura compuesta.

  \item \textbf{Columnas computadas.} Debe ser posible asociar una expresi\'on
    del DSL a una columna; el compilador genera la f\'ormula Excel
    correspondiente para cada fila de datos.

  \item \textbf{Referencias intrarregi\'on.} El lenguaje debe resolver
    referencias a columnas de la misma tabla (\dsl{col}) y a rangos de
    columnas de otra tabla en la misma regi\'on (\dsl{trange}, \dsl{tcol}).

  \item \textbf{Cuantificaci\'on existencial.} El lenguaje debe ofrecer una
    forma de expresar condiciones de existencia sobre filas de la misma tabla
    (\dsl{:all-rows}) o de otra tabla (\dsl{:rows-of}), especificando
    columnas de coincidencia y columnas de solapamiento.

  \item \textbf{Formato condicional persistente.} Las reglas de coloreado
    condicional deben traducirse a \texttt{FormulaRule} en el \xlsx{}, de
    modo que persistan y se reevaluen cuando el usuario edita el archivo.

  \item \textbf{Posicionamiento relativo.} Debe ser posible colocar
    expresiones fijas en posiciones calculadas como un ancla m\'as pasos
    direccionales, sin coordenadas absolutas codificadas.

  \item \textbf{Referencias interhoja.} El lenguaje debe soportar referencias
    din\'amicas a celdas de otras hojas mediante iteraci\'on sobre una lista
    de nombres de hojas (\dsl{collect-over}).

  \item \textbf{Ensamblaje declarativo.} El workbook completo debe poderse
    describir con formas del DSL (\dsl{libro}, \dsl{hoja}, \dsl{tabla})
    sin ning\'un conocimiento de las coordenadas del archivo de salida.

  \item \textbf{Independencia del backend.} El AST no debe contener
    ning\'un detalle del target de generaci\'on.  Agregar un nuevo tipo de
    expresi\'on no debe requerir modificar el DSL; solo a\~nadir un nuevo
    m\'etodo al backend.
\end{enumerate}

La tabla~\ref{tab:requisitos-casos} muestra c\'omo cada requisito se
manifiesta en los casos de uso analizados.

\begin{table}[h!]
\centering
\caption{Requisitos del lenguaje y casos de uso en que se manifiestan.}
\label{tab:requisitos-casos}
\renewcommand{\arraystretch}{1.2}
\begin{tabular}{p{5.2cm} ccc}
\toprule
\textbf{Requisito} & \textbf{Defensas} & \textbf{Facultad} & \textbf{TV} \\
\midrule
Declaraci\'on de tablas         & \checkmark & \checkmark & \checkmark \\
Columnas computadas              & \checkmark & \checkmark & \checkmark \\
Referencias intrarregi\'on       & \checkmark & \checkmark & --- \\
Cuantificaci\'on existencial     & \checkmark & --- & --- \\
Formato condicional persistente  & \checkmark & --- & --- \\
Posicionamiento relativo         & \checkmark & \checkmark & --- \\
Referencias interhoja            & --- & \checkmark & --- \\
Ensamblaje declarativo           & \checkmark & \checkmark & \checkmark \\
Independencia del backend        & \checkmark & \checkmark & \checkmark \\
\bottomrule
\end{tabular}
\end{table}
